\paragraph{A relative positive result over arbitrary bases of characteristic $p$.}

\begin{proposition}\label{prop:DE-E2-relative-Frobenius}
\status{Proved; independently audited by the parent and reduction-descent agents.}
Let $R$ be any scheme of characteristic $p$, and let $X\to R$ be locally of finite presentation.  Put $X^{(n)}=X\times_{R,F_R^n}R$.  Quasi-coherent crystals on $(X/R)_{\rm inf}$ are equivalent to systems of arbitrary quasi-coherent modules
\[
 M_n\in\operatorname{QCoh}(X^{(n)}),\qquad
 F_{X^{(n)}/R}^*M_{n+1}\simeq M_n.
\]
If every relative Frobenius $F_{X^{(n)}/R}$ is pure, restriction gives an equivalence between these crystals and ordinary stratified quasi-coherent $\OO_X$-modules.  No perfectness, reducedness, or Noetherian hypothesis on $R$ is required.
\end{proposition}
\begin{proof}
We verify that the construction of Lemma~\ref{lem:DE-crystal-Fdiv} is relative.  Locally on $X$ and $R$, choose a finitely presented closed immersion $X\hookrightarrow Y=\mathbf A_R^d$.  The iterated relative Frobenius $F_{Y/R}^n:Y\to Y^{(n)}$ is finite faithfully flat: in coordinates its module basis consists of the monomials with all exponents less than $p^n$.  More generally the same statement for smooth $Y/R$ is \cite[Lemma~50.19.5, Tag~0FW2]{StacksProject}.

Put $Y_n=Y\times_{Y^{(n)}}X^{(n)}$.  If $I=(f_1,\ldots,f_s)$ defines $X$ in $Y$, then $Y_n$ is defined by $(f_1^{p^n},\ldots,f_s^{p^n})$.  This remains true even though $F_R$ need not be flat: the ideal of a base-changed closed immersion is the image of the tensor product ideal, and relative Frobenius sends $f_i\otimes1$ to $f_i^{p^n}$.  Thus $X\hookrightarrow Y_n$ is nilpotent, and $Y_n\to X^{(n)}$ is finite faithfully flat.  Its double and triple fibre products are nilpotent thickenings of the diagonal copy of $X$.  Indeed, their defining difference ideals have finitely many generators killed by $p^n$th powers, in addition to the displayed nilpotent ideal.  A crystal therefore descends on this finite faithfully flat cover to a module $M_n$ on $X^{(n)}$.

The inclusion $Y_n\hookrightarrow Y_{n+1}$ and its maps on the double and triple products descend to $M_n\simeq F_{X^{(n)}/R}^*M_{n+1}$.  Conversely, for any nilpotent test $U\hookrightarrow T$ with ideal $J$, choose $n$ such that $J^{[p^n]}=0$.  In affine coordinates, the rule
\[
 (\OO_T/J)\otimes_{\OO_R,F_R^n}\OO_R
        \longrightarrow\OO_T,
       \qquad \overline b\otimes r\longmapsto b^{p^n}r
\]
is well defined and gives the relative Frobenius factor $\rho_n:T\to U^{(n)}$.  Set $\mathcal E_T=\rho_n^*(M_n|_{U^{(n)}})$.  The structural isomorphisms make this independent of $n$ and give the crystal transitions.

The inverse comparison is detected on
\[
                 Z=Y_n\times_{X^{(n)},\rho_n}T.
\]
It is finite faithfully flat over $T$ and is a nilpotent thickening of $U$.  For the latter assertion, modulo $J$ the extra difference ideal is finitely generated and killed elementwise by the $p^n$th power; adding $J$ preserves nilpotence.  This argument does not require $T$ to be Noetherian.  Ordinary faithfully flat module descent, including the double-product datum, proves that the constructions are inverse on all objects and morphisms.  Independence of the smooth embedding and gluing follow by the same canonical comparison.

Finally the relations $X\times_{X^{(n)}}X$ are cofinal among diagonal neighbourhoods: if $D$ is the diagonal ideal, their ideals are $D^{[p^n]}$, and the finite-generation bounds used in Proposition~\ref{prop:DE-E2-Fpure} apply over $R$ unchanged.  Under the stated purity hypotheses all iterated relative Frobenius maps are pure.  Effective module descent \cite[Theorem~35.4.22, Tag~08WE]{StacksProject} supplies $M_n$ from the stratification at every level; full faithfulness supplies the transition identifications.  This proves the second assertion.
\end{proof}

\begin{corollary}\label{cor:DE-E2-relative-node}
For every $\Fp$-algebra $R$, the flat finitely presented nodal family
\[
                    X=\Spec R[x,y]/(xy)\longrightarrow\Spec R
\]
satisfies the crystal--stratification equivalence for arbitrary quasi-coherent coefficients.
\end{corollary}
\begin{proof}
As an $R$-module its algebra is free on $1,x,x^2,\ldots,y,y^2,\ldots$, so the family is flat.  Every twist has the same displayed presentation.  The relative Frobenius $R[x_n,y_n]/(x_ny_n)\to R[x,y]/(xy)$ sends $x_n$ to $x^p$, $y_n$ to $y^p$, and fixes $R$.  It has an $R[x_n,y_n]/(x_ny_n)$-linear retraction: retain precisely the monomials whose exponent is divisible by $p$, divide that exponent by $p$, and keep their coefficients in $R$.  Linearity follows directly from $xy=0$.  Hence all relative Frobenius maps are split, in particular pure, and Proposition~\ref{prop:DE-E2-relative-Frobenius} applies.
\end{proof}
